\section{From Morse functions to Hamiltonian fixed points}
\label{dyn:section}
\label{sec:dynamics}

The dynamical construction turns a Morse function on
$X\times\CP^1$ into a Hamiltonian diffeomorphism of $X\times Y$ with
four times as many fixed points.  Here $Y$ will be a K\"ahler surface
with a Hamiltonian involution whose fixed locus consists of four
projective lines.  We give the construction of
\cite[Section~4]{OpenAICompanion}, including the argument that counts
all fixed points and checks their nondegeneracy in the ambient manifold.

We use the convention
\[
  \omega(X_f,\,\cdot\,)=df
\]
for the Hamiltonian vector field of a smooth function $f$.  A fixed
point $x$ of a diffeomorphism $\Phi$ is nondegenerate when
$d_x\Phi-I$ is invertible.  The following elementary estimate will
exclude additional fixed points after a small autonomous perturbation.
It is of the same type as Yorke's lower bound on periods for Lipschitz
vector fields~\cite{Yorke1969}; only a positive lower bound is needed here.

\begin{lemma}[A lower bound on nonconstant periods]
\label{dyn:short-period}
Let $V$ be a smooth vector field on a closed smooth manifold $P$, with
flow $\psi^s$.  There is a number $\tau>0$ such that
\[
  0<T<\tau,\qquad \psi^T(x)=x
  \quad\Longrightarrow\quad V(x)=0.
\]
\end{lemma}

\begin{proof}
Choose finitely many coordinate balls covering $P$, each with closure
contained in a larger coordinate ball in the same chart.  Bounds on
the coordinate speeds, together with the positive distances from
the smaller closed balls to the boundaries of the larger balls, give
$\eta>0$ such that a trajectory starting in a smaller ball stays in
its larger ball for time $\eta$.  Choose a common positive Lipschitz
constant $L$ for the coordinate expressions of $V$ on the larger
balls; their coordinate images are convex and their closures may be
taken compact in the charts.

Suppose $\psi^T(x)=x$ with $0<T<\eta$.  In a chart containing the whole
trajectory, write its coordinates as $y:[0,T]\to\R^d$, and write $v$
for the coordinate vector field.  Put
\[
  m=\max_{0\leq t\leq T}|v(y(t))|,
\]
and choose $t_0$ where the maximum is attained.  Since $y(T)=y(0)$,
\[
  \int_0^T v(y(t))\,dt=0,
  \qquad |y(t)-y(t_0)|\leq Tm.
\]
Consequently
\[
\begin{split}
  m
  &=\left|\frac1T\int_0^T
       \bigl(v(y(t_0))-v(y(t))\bigr)\,dt\right|\\
  &\leq\frac{L}{T}\int_0^T |y(t_0)-y(t)|\,dt
   \leq LTm.
\end{split}
\]
For $T<\tau:=\min\{\eta,1/L\}$ this forces $m=0$, proving the claim.
\end{proof}

The next proposition applies the estimate to an involution.  Its
fixed components are symplectic submanifolds.  An invariant extension
of Morse functions on those components supplies the perturbation;
squaring the perturbed map will identify every one of its fixed points.

\begin{proposition}[Perturbing a Hamiltonian involution]
\label{dyn:involution}
Let $(P,\omega)$ be a closed symplectic manifold and let $g:P\to P$
be a Hamiltonian diffeomorphism with $g^2=\mathrm{id}_P$.  Write
$\Fix(g)=F_1\sqcup\cdots\sqcup F_r$ for its connected fixed
components.  For each $j$, let $h_j:F_j\to\R$ be a smooth Morse
function.  There is a smooth $g$-invariant function $f:P\to\R$
restricting to $h_j$ on $F_j$ such that, for all sufficiently small
$\epsilon>0$, its Hamiltonian flow $\psi^s$ satisfies
\begin{equation}
\label{dyn:fixed-set}
  \Fix(\psi^\epsilon\circ g)
    =\bigsqcup_{j=1}^r\Crit(h_j).
\end{equation}
Every fixed point in this equality is nondegenerate.  Moreover,
$\psi^\epsilon\circ g$ is the time-one map of a smooth one-periodic
Hamiltonian.  If $P$ is simply connected, all fixed-point loops along
the resulting Hamiltonian path are contractible.
\end{proposition}

\begin{proof}
Average a Riemannian metric over $\{\mathrm{id}_P,g\}$.  At a fixed
point $x$, the resulting exponential chart intertwines $g$ with
$d_xg$.  Thus the fixed locus is locally a smooth submanifold,
whose tangent space at $x$ is
\[
  E_x^+=\ker(d_xg-I).
\]
This local description and compactness imply that it has finitely
many connected components, each a smooth closed submanifold.
The involution gives the
decomposition
\[
  T_xP=E_x^+\oplus E_x^- ,
  \qquad E_x^-=\ker(d_xg+I).
\]
For $u\in E_x^+$ and $v\in E_x^-$, symplecticity gives
\[
  \omega_x(u,v)
  =\omega_x(d_xg(u),d_xg(v))=-\omega_x(u,v).
\]
Thus the summands are symplectically orthogonal.  Since $\omega_x$
is nondegenerate, its restriction to each summand is nondegenerate
as well.  In particular, $F_j$ is symplectic and $E_x^-$ is its
symplectic normal space.

Using disjoint tubular neighborhoods of the $F_j$ and smooth cutoff
functions, extend the $h_j$ to a smooth function $\widetilde f$ on $P$.
Set
\[
  f=\tfrac12(\widetilde f+\widetilde f\circ g).
\]
This function is $g$-invariant and restricts to $h_j$ on $F_j$.
Invariance and symplecticity imply $g_*X_f=X_f$, so $g$ commutes with
$\psi^s$.  At a point $x\in F_j$, invariance also implies
$d_xf(v)=0$ for $v\in E_x^-$.  Orthogonality of the two summands then
shows that $X_f$ is tangent to $F_j$, where it is the Hamiltonian
vector field of $h_j$ for $\omega|_{F_j}$.  Hence
\begin{equation}
\label{dyn:equilibria}
  x\in F_j:\qquad X_f(x)=0
  \quad\Longleftrightarrow\quad d_xh_j=0.
\end{equation}

Put $\Phi_\epsilon=\psi^\epsilon\circ g$.  The commutation relation
and $g^2=\mathrm{id}_P$ give
\[
  \Phi_\epsilon^2=\psi^{2\epsilon}.
\]
Let $\tau$ be given by Lemma~\ref{dyn:short-period} for $X_f$, and
take $2\epsilon<\tau$.  Any fixed point $x$ of $\Phi_\epsilon$
satisfies $\psi^{2\epsilon}(x)=x$, so the lemma gives $X_f(x)=0$.
The flow therefore fixes $x$ at every time, and
$\psi^\epsilon(g(x))=x$ implies $g(x)=x$.  Equation~\eqref{dyn:equilibria}
now places $x$ among the prescribed critical points.  Conversely,
each prescribed critical point is fixed by both $g$ and $\psi^\epsilon$.
This proves~\eqref{dyn:fixed-set}.

We next check nondegeneracy in both the tangent and the normal
directions.  At a point $x$ in the finite set~\eqref{dyn:fixed-set},
let $A_x=d_xX_f$ be the linearization at the equilibrium.  Then
\[
  d_x\psi^\epsilon=\exp(\epsilon A_x).
\]
Equivariance implies that $A_x$ commutes with $d_xg$, so it preserves
$E_x^+$ and $E_x^-$.  Denote its restrictions by $A_x^+$ and $A_x^-$.
On $E_x^+=T_xF_j$, differentiating the Hamiltonian equation at the
critical point gives
\[
  \omega_x(A_x^+u,v)=\operatorname{Hess}_x h_j(u,v).
\]
The Morse condition therefore makes $A_x^+$ invertible.  Since
\[
  \frac{\exp(\epsilon A_x^+)-I}{\epsilon}
       \longrightarrow A_x^+
       \quad\text{as }\epsilon\longrightarrow0,
\]
the tangent block of $d_x\Phi_\epsilon-I$ is invertible for all
sufficiently small positive $\epsilon$.  The normal block is
\[
  -\exp(\epsilon A_x^-)-I\longrightarrow-2I,
\]
and is likewise invertible.  There are finitely many prescribed
critical points, so a single positive upper bound for $\epsilon$
ensures nondegeneracy at all of them as well as the fixed-set equality.

Finally, choose a smooth Hamiltonian $K_s$, $0\leq s\leq1$, whose
time-one map is $g$.  Choose nonnegative smooth functions $a,b$ with
integral one and supports contained in two successive disjoint
subintervals of $(0,1)$.  Define
\[
  A(t)=\int_0^t a(u)\,du,\qquad
  H_\epsilon(t,x)=a(t)K_{A(t)}(x)+\epsilon b(t)f(x),
  \qquad 0\leq t\leq1.
\]
During the support of $a$, the flow traverses the path of $K_s$ from
the identity to $g$.  During the later support of $b$, it follows
the flow of $f$ for total time $\epsilon$.  Thus its time-one map
is $\psi^\epsilon\circ g$.  Since $H_\epsilon$ vanishes near both
endpoints of $[0,1]$, its periodic extension is smooth.  If $P$ is
simply connected, every fixed-point loop of this specified path is
contractible.
\end{proof}

We now construct the surface that gives four copies of the Morse
problem.  Blowing up each isolated fixed point of a half-turn
replaces it by a fixed projective line.

\begin{proposition}[The surface with four fixed curves]
\label{dyn:surface}
Let $Y$ be the complex blow-up of $\CP^1\times\CP^1$ at
\[
  \{0,\infty\}\times\{0,\infty\}.
\]
There is a circle-invariant K\"ahler form $\omega_Y$ for which the
lifted diagonal circle action is Hamiltonian.  Its half-turn $g_Y$
has fixed locus precisely the four exceptional curves
$E_1,\ldots,E_4$, each isomorphic to $\CP^1$.  Along these curves,
$d g_Y$ is the identity on the tangent bundle and minus the identity
on the symplectic normal bundle.  Moreover, $Y$ is simply connected
and its integral homology is
\begin{equation}
\label{dyn:Y-homology}
  H_i(Y;\Z)\cong
  \begin{cases}
    \Z,&i=0,4,\\
    \Z^6,&i=2,\\
    0,&\text{otherwise}.
  \end{cases}
\end{equation}
\end{proposition}

\begin{proof}
In affine coordinates, the diagonal circle action on
$\CP^1\times\CP^1$ is
\[
  t\cdot(z,w)=(e^{2\pi it}z,e^{2\pi it}w),
  \qquad t\in\R/\Z.
\]
It preserves the four centers and therefore lifts holomorphically
to $Y$.  Proposition~\ref{geo:blowup} shows that $Y$ is K\"ahler and
simply connected.  Average a K\"ahler form over the lifted circle
action to obtain an invariant K\"ahler form $\omega_Y$.  If $V$ is
the infinitesimal generator, Cartan's formula gives
\[
  d\bigl(\omega_Y(V,\,\cdot\,)\bigr)
     =\mathcal L_V\omega_Y=0.
\]
As $H^1(Y;\R)=0$, there is a smooth function $k:Y\to\R$ with
$dk=\omega_Y(V,\,\cdot\,)$.  Thus $V=X_k$, and the autonomous
Hamiltonian $k/2$ has time-one map $g_Y$.

The downstairs half-turn fixes exactly the four centers.  Using
the coordinate $z$ near $0$ and $1/z$ near $\infty$ in each factor,
the circle weights at the centers are
\[
  (1,1),\quad(1,-1),\quad(-1,1),\quad(-1,-1).
\]
For weights $(\alpha,\beta)$ the induced action on the exceptional
line is
\[
  [z_1:z_2]\longmapsto
  [e^{2\pi i\alpha t}z_1:e^{2\pi i\beta t}z_2].
\]
When the signs differ this action has nonzero relative weight
$\beta-\alpha=\pm2$; in all four cases its value at $t=1/2$ is the
identity.  To see the action also in the normal direction, use the
local blow-up model
\[
  \{(\ell,v):\ell\in\CP^1,\ v\in\ell\subset\C^2\}.
\]
The downstairs half-turn is $v\mapsto-v$, so its lift is
\[
  (\ell,v)\longmapsto(\ell,-v).
\]
It fixes exactly the zero section, with derivative $+I$ on its
tangent space and $-I$ on its normal line.  The symplectic normal
space is the negative eigenspace by the orthogonality argument in
Proposition~\ref{dyn:involution}.  Away from the exceptional curves,
the blow-down is an equivariant isomorphism.  There are therefore
no additional fixed points.

Finally, each point blown up has complex codimension two, so
\eqref{geo:blowup-homology} adds one free generator in degree two.
The integral homology of $\CP^1\times\CP^1$ is free with ranks
$(1,2,1)$ in degrees $(0,2,4)$, respectively.  Four blow-ups give
\eqref{dyn:Y-homology}.
\end{proof}

\begin{corollary}[Product realization]
\label{dyn:product}
Let $X$ be a connected closed simply connected K\"ahler manifold,
let $\omega_X$ be any K\"ahler form on $X$, and let $(Y,\omega_Y)$ be as in
Proposition~\ref{dyn:surface}.  Give $M=X\times Y$ the product form
$\omega_X\oplus\omega_Y$ and put $C=X\times\CP^1$.
For every smooth Morse function $h:C\to\R$, there is a smooth
one-periodic Hamiltonian $H$ on $M$ whose time-one map has only
nondegenerate fixed points and satisfies
\begin{equation}
\label{dyn:product-count}
  \#\Fix(\phi_H^1)
   =\#\Fix_0(\phi_H^1;H)
   =4\,\#\Crit(h).
\end{equation}
\end{corollary}

\begin{proof}
The pullback to $M$ of $k/2$ generates the Hamiltonian involution
$g=\mathrm{id}_X\times g_Y$ for the given product form.  Its fixed
components are exactly $X\times E_1,\ldots,X\times E_4$.  Choose
isomorphisms $E_j\cong\CP^1$ and use them to pull $h$ back to Morse
functions $h_j$ on these four components.  Each $h_j$ has
$\#\Crit(h)$ critical points.  Proposition~\ref{dyn:involution}
applies to their induced symplectic forms and produces exactly
the fixed points counted in~\eqref{dyn:product-count}, all
nondegenerate.  Finally, $M$ is simply connected, so every
fixed-point loop for the constructed path is contractible.
\end{proof}
